%preamble{
\usepackage{beamer-JA}
\usepackage{tikz,calc}
\usetikzlibrary{calc,positioning,arrows,arrows.meta,shapes, shapes.geometric,fit}
\usepackage{adjustbox}
\usepackage{wasysym} %for symbols
% \usepackage{venndiagram}
\usepackage{caption}
%\usepackage{booktabs}
%\usepackage{caption,subcaption}
\usepackage{amssymb}
\usepackage{mathtools}

\DeclareMathAlphabet{\mathcal}{OMS}{cmsy}{m}{n}
\SetMathAlphabet{\mathcal}{bold}{OMS}{cmsy}{b}{n} %for big-O
% \setbeamertemplate{footline}[page number] %for coding
% \setbeamertemplate{footline}[totalframenumber]
% \beamerdefaultoverlayspecification{<+->}
\newcounter{num}
\setcounter{num}{1}
\newcounter{numTemp}
\setcounter{numTemp}{1}
%%%%Chapter and Section numbers%%%%% 
% \newcommand{\chNum}{}
% \newcommand{\secNum}{}
\newcommand{\secTitle}{C964 section title}
\date{}
%\date{\formatdate{29}{9}{2015}}

%tikz settings{
%}
\definecolor{blueWGU}{RGB}{0,48,87}
\definecolor{lightBlueWGU}{RGB}{73,134,160}

\definecolor{uablue}{RGB}{0,61,100}
\colorlet{uablue100}{uablue}
\colorlet{uablue75} {uablue!75!white}
\colorlet{uablue50} {uablue!50!white}
\colorlet{uablue25} {uablue!25!white}
\colorlet{uablue10} {uablue!10!white}
\colorlet{uablue5}  {uablue!5!white}

\definecolor{uared}{RGB}{126,0,47}
\colorlet{uared100}{uared}
\colorlet{uared75} {uared!75!white}
\colorlet{uared50} {uared!50!white}
\colorlet{uared25} {uared!25!white}
\colorlet{uared10} {uared!10!white}
\colorlet{uared5}  {uared!5!white}

\tikzset{onslide/.code args={<#1>#2}{%
  \only<#1>{\pgfkeysalso{#2}} % \pgfkeysalso doesn't change the path
}}
\tikzset{alt/.code args={<#1>#2#3}{%
  \alt<#1>{\pgfkeysalso{#2}}{\pgfkeysalso{#3}} % \pgfkeysalso doesn't change the path
}}
\tikzset{temporal/.code args={<#1>#2#3#4}{%
  \temporal<#1>{\pgfkeysalso{#2}}{\pgfkeysalso{#3}}{\pgfkeysalso{#4}} % \pgfkeysalso doesn't change the path
}}
%for fancy arrows
\tikzset{
node distance = 9em and 4em,
%sloped,
   box/.style = {%
    shape=rectangle,
    rounded corners,
    draw=blue!40,
    fill=blue!15,
    align=center,
    font=\fontsize{12}{12}\selectfont},
    double -latex/.style args={#1 colored by #2 and #3}{
        -latex,line width=#1,#2,
        postaction={draw,-latex,#3,line width=(#1)/3,shorten <=(#1)/4,shorten >=4.5*(#1)/3},
    },
    arrow/.style args= {#1}{%
        draw=black,
        line width=2mm,% <-- select desired width
        -{Triangle[length=3.25mm]},
        shorten >=1mm, shorten <=1mm,
        font=\fontsize{8}{8}\selectfont,
        postaction={draw=#1!35,line width=1.75mm,-{Triangle[length=2.75mm]},shorten >=1.25mm,shorten <=1.25mm}
        },
    arrow/.default={blue} %arrow2/.default={blue}{yellow}
}

\tikzstyle{block} = [rectangle, draw, fill=uablue25,
    text width=4.5em, text centered, rounded corners, minimum height=3.25em, line width=1pt ]
\tikzstyle{block2} = [rectangle, draw, fill=uablue25,
    text width=3.0em, text centered, rounded corners, minimum height=4.0em, line width=1pt ]
\tikzstyle{block3} = [rectangle, rounded corners, line width = 1pt, draw, fill=uablue25, text centered, minimum width=6em, minimum height = 2em, inner sep =0em]

\tikzstyle{mysquare} =[rectangle, draw, text width=4em, text centered, rounded corners, minimum height=2cm, minimum width=2cm, line width=2pt,fill=red!50!white]
\tikzstyle{mydiamond} =[diamond, draw, text width=4em, text centered, rounded corners, minimum height=2cm, minimum width=2cm, line width=2pt,fill=blue!50!white]
\tikzstyle{mycircle} =[circle, draw, text width=4em, rounded corners, line width=2pt]
%\tikzstyle{subsetset} = [thick,aspect=0.7,minimum height=1.7cm,minimum
%                            width=1.5cm, align = center]
\tikzstyle{cylinderl} = [
    minimum width=1.5 cm, minimum height=1.75 cm,
    inner sep=0pt,
    align=center]

% cylinder node
\tikzset{
  pics/myCylinder/.style args = {#1}{ %#1/#2/#3.. for multiple args
    code = {
%    \pgfsize{\mywidth}{\myheight} %can use \mywidth to direcly get box value
    \node[name=origin] at (0,0) { }; %places one box corner at 0,0
    \useasboundingbox (origin);
    \def\cTop{.2}; %1cm = 28.3464567 pt
%    \coordinate (UR) at (\mywidth/2,\myheight/2-\cTop);
    \coordinate[below=\cTop of current bounding box.north east] (UR);
    \coordinate[above=\cTop of current bounding box.south east] (LR);
    \coordinate[above=\cTop of current bounding box.south west] (LL); %correct
    \coordinate[below=\cTop of current bounding box.north west] (UL);
    \coordinate[below=\cTop of current bounding box.north] (top);

%    \draw[thick, green] let \p1 = (UR), \p2 = (UL) in (\x1-\x2,0) -- (LR) -- (LL) -- (UL) -- cycle;
    \draw[fill = #1,draw] let \p1 =(UR), \p2 = (UL) in (LL) arc[start angle=180, end angle = 360, radius =\x1/2-\x2/2, y radius = \cTop cm] -- (UR) -- (UL) -- (LL) -- cycle;
    \draw [fill =#1!75!black] let \p1 =(UR), \p2 = (UL) in (top) ellipse (\x1/2-\x2/2 and \cTop cm); %top
    \draw (LL)--(UL); %left
    \draw (UR)--(LR); %right
    \draw let \p1 =(UR), \p2 = (UL) in (LL) arc (180:360:\x1/2-\x2/2 and \cTop cm); %bottom front base
    \node[font=\Large] at (0,-\cTop/2) {\tikzpictext}; %[block] to see outline for placement
    }
  }
}
\tikzstyle{line} = [line width=1pt, -triangle 45]
\tikzset{
  mycustomnode/.style={
    draw,
    single arrow,
    single arrow head extend=0,
    shape border uses incircle,
    shape border rotate=-90,
  }
}
% \tikzstyle{alert} = [text=uared100, fill=uared25, draw=uared100]
\tikzstyle{dim} = [text=uablue25, fill=uablue5, draw=uablue25]
\tikzstyle{highlight} = [fill=green!50!white, line width =2pt]
\tikzstyle{off} = [text=white, fill=white, draw=white,line width =2pt]
\tikzstyle{dim} = [opacity=.25]
\tikzstyle{myframe} = [line width=4pt, draw=green!65!blue,inner sep=.5em,rounded corners]


% %%for counting slides and frames when debugging
% \newcounter{slidenumber}
% \defbeamertemplate*{footline}{infolines theme frame plus slide}{
%     \setcounter{slidenumber}{\insertpagenumber}%
%     \addtocounter{slidenumber}{-\insertframestartpage}%
%     \addtocounter{slidenumber}{1}%
%     \leavevmode%
%     \hbox{%
%         \begin{beamercolorbox}[wd=\paperwidth,ht=2.25ex,dp=1ex,center]{date in head/foot}%
%           \ifthenelse{\insertframeendpage=\insertframestartpage}{
%             \insertframenumber/ \inserttotalframenumber\hspace*{2ex}}{
%             \insertframenumber(\arabic{slidenumber}{})/\inserttotalframenumber\hspace*{2ex}}
%         \end{beamercolorbox}}%
%         \vskip0pt%
%     }
% \setbeamertemplate{footline}[infolines theme frame plus slide]
\usepackage{xcolor,soul}
% \definecolor{lightblue}{rgb}{.90,.95,1}
\sethlcolor{green}
\renewcommand<>{\hl}[1]{\only#2{\beameroriginal{\hl}}{#1}}

%%higlighting tex wiht overlays
\makeatletter
\newcommand\SoulColor{%
  \let\set@color\beamerorig@set@color
  \let\reset@color\beamerorig@reset@color}
\makeatother
\SoulColor
%}